% Source r6 complements, record catalogue, and formal details.
\section{Complements and further questions}\label{sec:open}

Problem~\#243 is open.  The negative case has narrowed, and it is worth being
precise about what is left.

Each exclusion holds only in its stated regime.  Theorems~\ref{res:constant}
and~\ref{res:periodic} exclude the constant and the stated periodic magnitudes
on an all-negative tail, and do not exclude those patterns merely along the
negative indices of a mixed-sign tail; Theorems~\ref{res:bounded}
and~\ref{res:mass} exclude a bounded negative part and finite normalised
negative mass under normalised vanishing.  Together with absorption and descent
they fix the profile of any counterexample.

\begin{proposition}[frontier profile for a counterexample]\label{res:frontier}
For the canonical integer state attached to any counterexample to
Problem~\ref{res:problem},
\[
 E_n\ne0\quad\hbox{eventually},\qquad
 \frac{|E_n|}{C_n}\longrightarrow0,
\]
and
\[
 \limsup_{\substack{n\to\infty\\E_n<0}}(-E_n)=\infty,
 \qquad
 \sum_{n=0}^{\infty}\frac{(-E_n)_+}{C_n}=\infty .
\]
In particular, negative indices occur infinitely often.  That last clause is
not new: it is the contrapositive of Koizumi's Proposition~1(2), which he
attributes to Badea, and it is what the present theorem strengthens by
replacing eventual nonnegativity with eventual boundedness of the negative
part.  Thus the remaining
regime consists of globally coherent, unbounded negative excursions with
divergent normalised negative mass.
\end{proposition}

\begin{proof}
Koizumi's canonical-tail transfer gives strict centring and normalised
vanishing after a finite shift; the strict-centring hypothesis of
Theorem~\ref{res:bounded} also follows from normalised vanishing with $K=1$.
Absorption then makes $E_n\ne0$ eventually,
since eventual zero would give the Sylvester recurrence.  Descent excludes an
eventually nonnegative state.  Theorem~\ref{res:bounded} excludes a bounded
negative part, and Theorem~\ref{res:mass} excludes finite normalised negative
mass.  These statements are unchanged by deleting a finite prefix.
\end{proof}

Put
\[
 \delta_n:=\frac{(-E_n)_+}{C_n}.
\]
Every argument in Sections~\ref{sec:constant}--\ref{sec:bounded} uses a
finiteness hypothesis: a fixed set of prime divisors, a fixed period, or a
fixed bound on the negative part.  Theorem~\ref{res:mass} uses
$\sum_n\delta_n<\infty$.  Proposition~\ref{res:frontier} records that a
counterexample has none of them; it is a frontier statement, not an additional
problem equivalent to the original one.  Stated at the level of state orbits,
the surviving obstruction is the following.

\begin{problem}[unbounded divergent-mass negative excursions]\label{res:excursions}
Exclude, or construct, an exact natural orbit $(a,D,C)$ with $a_n>1$ and
$C_n>0$, whose multipliers satisfy $\lim_n a_{n+1}/a_n^{2}=1$ and whose centred
state satisfies normalised vanishing, and which is
negative infinitely often with magnitudes unbounded along that cofinal set and
\[
 \sum_n\frac{(-E_n)_+}{C_n}=\infty .
\]
An excursion of this kind has none of the finiteness hypotheses just listed.
By Proposition~\ref{res:frontier} the canonical state of a counterexample is
such an orbit after deletion of a finite prefix, so an exclusion would settle
Problem~\#243.  Finite states and scalar error models need not realise a
reciprocal series.  An infinite positive exact state with the displayed
recurrences and normalised vanishing does: telescoping gives
$\sum_{k\ge n}1/a_k=C_n/D_n$ and $a_{n+1}/a_n^2\to1$, so after deleting a
finite prefix such a nonterminal orbit is a counterexample to
Problem~\ref{res:problem}.
\end{problem}

\begin{remark}
The hypotheses in Problem~\ref{res:excursions} beyond the state recurrence are
not decorative: the recurrence alone admits unbounded divergent-mass
excursions.  Put
\[
 C_n=2^{n},\qquad b_0=2,\qquad b_{n+1}=\tfrac12 b_n(b_n+2),\qquad
 a_n=b_n+2,\qquad D_n=b_nC_n ,
\]
so that $(b_n)=2,4,12,84,3612,\ldots$ and $(a_n)=4,6,14,86,3614,\ldots$.  Every
$b_n$ is a positive even integer, since $b_n=2k$ gives $b_{n+1}=2k(k+1)$, so all
four sequences are $\N$-valued.  The orbit is exact:
$C_{n+1}+D_n=2^{n}(2+b_n)=a_nC_n$ and
$D_{n+1}=b_{n+1}2^{n+1}=b_n(b_n+2)2^{n}=a_nD_n$.  Its centred state is
\[
 E_n=D_n-(a_n-1)C_n=b_n2^{n}-(b_n+1)2^{n}=-C_n ,
\]
so the negative magnitudes $2^{n}$ are unbounded and $\sum_n(-E_n)_+/C_n$
diverges.  What fails is the centring: $|E_n|=C_n$ at every index, so strict
centring fails at the boundary and normalised vanishing fails outright.  The
multipliers satisfy $a_{n+1}=\tfrac12(a_n^{2}-2a_n+4)$, so
$a_{n+1}/a_n^{2}\to\tfrac12$ and the growth hypothesis fails as well.  The
example therefore says nothing about Problem~\ref{res:problem}, and it does not
isolate the centring hypotheses: $E_n=-C_n$ at every index forces that
multiplier recurrence, so growth fails with them.  It is recorded only to show
that the state recurrence alone is not enough.
\end{remark}

\subsection{Global height and old-factor overlap}

The main arithmetic gap is global.  Define the cumulative digit LCM and its
overlap quotient by
\[
 L_0=D_0,\qquad L_{n+1}=\lcm(L_n,a_n),\qquad M_n=\frac{D_n}{L_n}.
\]
Since $D_n=D_0\prod_{j<n}a_j$, the quotient is integral and
$M_nL_n=D_n$.  It records the multiplicity lost when the full product scale
is compressed to an LCM.

\subsection{Repair entropy across a recovery}

The cumulative quotient does not by itself distinguish one large deletion
from many small ones.  The reduced recurrence does.  Let $u,h,c:\N\to\N$ and
$e:\N\to\Z$ satisfy
\[
 h_nu_{n+1}=u_n-e_n,
 \qquad u_n>0,
 \qquad h_n>0.
\]
Fix $K,r,L\in\N$ with $K,L>0$, assume
\[
 K|e_{r+i}|<u_{r+i}\quad(0\le i<L),
 \qquad u_r\le u_{r+L},
\]
and let $R$ be a finite family of moduli.  Suppose every $m_q$, $q\in R$,
divides the deletion product $\prod_{r\le n<r+L}c_n$, and that
$c_n^2\mid h_n$ throughout the interval.  Then
\begin{equation}\label{eq:repair-entropy}
 K^L\lcm(m_q:q\in R)^2<(K+1)^L.
\end{equation}
This is the
\lword{Erdos243/RepairEntropy.lean}{48}
      {repairedFamily_recovery_energy_divisionFree}
      {repair-entropy conservation law}.

The square in~\eqref{eq:repair-entropy} is the key point.  Exact denominator
reduction charges $c_n^2$ to the scale growth, so a repaired family whose LCM
is at least $2^{|R|}$ satisfies
\[
 K^L4^{|R|}<(K+1)^L.
\]
The formal consequence is
\lword{Erdos243/RepairEntropy.lean}{78}
      {repaired_card_bound_of_independent}
      {the independent-repair cardinality bound}.
It rules out positive-density deletion of independent old moduli inside
recoveries with a fixed relative-error budget.

There is also an endpoint that does not mention a chosen modulus family.  If
$K|e_n|<u_n$ eventually for every $K$, then for each fixed $L>0$ there is an
$N$ such that every recovery $u_r\le u_{r+L}$ with $r\ge N$ satisfies
\[
 \prod_{0\le i<L}h_{r+i}=1.
\]
Thus a surviving nontrivial reset must have recovery length tending to
infinity.  Lean checks this as
\lword{Erdos243/RepairEntropy.lean}{107}
      {eventually_recoveryPayment_eq_one_of_fixedLength}
      {eventual disappearance of fixed-length reset payments}.
The theorem does not bound a recovery length that varies with $r$, nor does it
supply a lower bound for the LCM of the moduli repaired in such an interval.
Those are the missing inputs needed to turn~\eqref{eq:repair-entropy} into a
global contradiction.

\begin{problem}[global overlap-height growth]\label{res:lcmheight}
For every nonterminal canonical orbit satisfying Problem~\ref{res:problem},
must
\[
 \limsup_{n\to\infty}\frac{\log M_n}{n}>0?
\]
Equivalently, must there be a $K\ge1$ for which
\[
 2^n\le M_n^K
\]
at infinitely many indices?
\end{problem}

\noindent The two displayed formulations are equivalent up to changing the
positive constant; the second is not a weaker target. On the canonical orbit,
Section~\ref{sec:lcmrecords} supplies $M_n\mid C_n$ and $M_n\le C_n$.
Therefore a positive answer would contradict the subexponential tail-height
budget and prove Problem~\#243. The missing input is the proposed lower
bound on overlap growth.

There is also a more local-looking question whose content is nevertheless the
entire prefix.  Put $A_n=\prod_{j<n}a_j$, so $D_n=D_0A_n$.  From
\[
 D_0A_n=(a_n-1)C_n+E_n
\]
one immediately obtains
\[
 \gcd(A_n,a_n-1)\mid E_n.
\]
The checked
\lword{Erdos243/ReciprocalTailRigidity.lean}{2164}{tailState_subexponential_of_normalizedVanishes}{tail-height estimate}
and normalised vanishing make $|E_n|$ subexponential in $n$.

\begin{problem}[old-factor overlap]\label{res:prefixgcd}
In every nonterminal canonical orbit satisfying Problem~\ref{res:problem}, is
\[
 \limsup_{n\to\infty}
 \frac{\log\gcd(A_n,a_n-1)}{n}>0?
\]
Equivalently, do there exist $\eta>0$ and infinitely many $n$ such that
$\gcd(A_n,a_n-1)\ge e^{\eta n}$?
\end{problem}

\noindent A positive answer contradicts the displayed divisibility and the
subexponential error budget.  Arbitrarily long locally admissible blocks do
not answer this question: the gcd uses the complete prefix.

\subsection{The direct analytic question}

\begin{problem}[summability from rationality]\label{res:masshyp}
Let $(a_n)$ satisfy Problem~\ref{res:problem}, and let $(D_n,C_n,E_n)$ be its
canonical integer state.  Must
\[
 \sum_{n=0}^{\infty}\delta_n
 =\sum_{n=0}^{\infty}\frac{(-E_n)_+}{C_n}<\infty ?
\]
\end{problem}

\noindent An affirmative answer closes Problem~\#243 by
Theorem~\ref{res:mass}.  Since $\delta_n\to0$, convergence is equivalently
expressible as boundedness of the partial products
$\prod_{n<N}(1+\delta_n)$; this is the same criterion in multiplicative form.
A negative answer requires the full infinite hypotheses.  A finite admissible
path or a scalar model is insufficient.  An infinite positive exact $C/D$
orbit with $a_n\ge2$ and $E_n/C_n\to0$, however, already realises a rational
reciprocal tail by telescoping; it should not be dismissed as merely local.
By Theorem~\ref{res:weightedrecord}, it is enough to establish the finiteness
of~\eqref{eq:weightedgrowth} for one decreasing weight with divergent
integral and one fixed baseline. This replaces the all-index mass by a
discounted series supported only at LCM records.

\begin{remark}
One further question is not left open here, since it is answered
in~\cite{koizumi2025}: whether normalised vanishing follows from
$a_{n+1}/a_n^{2}\to1$ and rationality.  It does.  After
deleting a finite prefix, Koizumi's Corollary~3
~\cite[p.~9]{koizumi2025} makes the sequence the pseudo-greedy expansion of its
own reciprocal sum and gives $\varepsilon_n\to0$ under exactly the hypotheses of
Problem~\ref{res:problem}; rationality is not needed for that step, only
summability of $\sum1/a_n$.  Since $\varepsilon_n=E_n/C_n$ under the intended
reading of Section~\ref{sec:state}, this is normalised vanishing.  Strict
centring comes from the same source and is stronger than what
Theorem~\ref{res:bounded} assumes: for a rational reciprocal sum, the range
$-c_n/2\le e_n<c_n/2$ of his Lemma~4(2)
~\cite[p.~11]{koizumi2025} gives $|E_n|<C_n$ at every index, which is stronger
than eventual strict centring.  The eventual form also follows directly from
normalised vanishing with $K=1$.  What remains unsupplied is hypothesis~(5), the bound
on the negative part, which is why the theorem is still conditional.
\end{remark}

\subsection{A barrier extension}

\begin{problem}[growth of the moduli in a variable-rise barrier]\label{res:variablerise}
Let $m_j$ be pairwise-coprime old moduli and let $u_n$ be an unbounded
integer sequence avoiding them as in Theorem~\ref{res:barrier}.
Which joint conditions on $(u_{n+1}-u_n)_+$ and on the growth and activation
indices of the $m_j$ force a contradiction?  In particular, determine the
consequences of the double-logarithmic multiplier growth supplied by a
canonical reciprocal-tail orbit.
\end{problem}

\noindent For each unbounded nondecreasing $\omega$, there is a prime family
and a coprime-avoidance sequence with unbounded gaps $o(\omega(u_n))$.
The family and sequence depend on $\omega$.  Thus no such envelope works
uniformly over arbitrary modulus families; this is not a reciprocal-tail
counterexample.

\subsection*{Formalisation targets}

\paragraph{Erd\H{o}s--Straus.}
Formalise Theorem~3 of~\cite{erdosstraus1964}, with its prefix-LCM quotient
and correctly indexed growth factor, under the published analytic hypotheses.

\paragraph{Duverney.}
Formalise Corollary~3.2 of~\cite{duverney2001}, including its signed
numerators and its stated defect-series convergence hypothesis, and then
recover the all-positive specialisation used here.  Under a merely conditional reading of
the defect series, convergence of $\sum(r_n-1)$ does not force
$\inf\prod r_n>0$; a positive-term absolutely summable specialisation is the
safe formalisation target.  A self-contained signed specialisation assuming
also $\sum\eta_n^2<\infty$ makes the product argument elementary; its algebraic
recurrences and cleared telescoping identity are in
\texttt{SignedDuverneyAlgebra.lean}.  That trap does not refute the published
theorem.  These are useful verification projects, but neither is an additional
open mathematical problem.

\pdfbookmark[1]{Statements and declarations}{statements-declarations}
\subsection*{Statements and declarations}

\paragraph{Artefact and data availability.} The
\href{\sourceurl}{pinned formal-source revision} contains the Lean sources, the
fixed toolchain, and the library manifest used in the verification.  The linked Lean declarations provide formal evidence for their exact
statements.  Ordinary theorems in this manuscript retain their displayed
proofs; they are not certified merely by nearby source links.

Lean checks each proof term against the fixed library
version, and the sources linked here contain no proof placeholders and no
project-defined axioms; Lean does not authorise the exposition, the citation
choices, or the interpretation, for which the author remains responsible.

\paragraph{Funding and competing interests.} This work received no external
funding.  The author declares no competing interests.

\paragraph{Acknowledgements.} The problem numbering and status follow the
Erd\H{o}s Problems catalogue maintained by Thomas Bloom~\cite{erdosproblems}.
